\section{Limitations and nonclaims}\label{sec:limitations-nonclaims}


This section collects the scope restrictions of
Theorem~\ref{thm:scoped-exact-six} and the global nonclaims
preserved throughout.

\subsection{Scope of the terminal theorem}\label{subsec:terminal-scope}

Theorem~\ref{thm:scoped-exact-six} establishes
\(\ScopedExactSix{\FATCD_{\mathrm{RH}}}\); its six-channel decomposition
clauses hold for all \(D\in\FATCD{}\), the finite audited typed closure
description domain of Definition~\ref{def:fatcd}. The theorem asserts:
\begin{itemize}[topsep=2pt,itemsep=1pt]
  \item existence of a decomposition/exhaustion record \(\Dec{D}\)
    with six role channels and one status per channel;
  \item residual-feature exhaustion under the upper-bound classification,
    conditional on the recognition hypothesis
    (Definition~\ref{def:recognition-hypothesis}) and within the covered scope,
    so that no seventh irreducible role is required there (the three principal
    threat classes are closed without the hypothesis, in the recognized
    forms of
    Propositions~\ref{prop:probability-closure}--\ref{prop:agency-closure});
  \item activation of any \(\pi_i(D)\) only under the threshold,
    witness-schema, collapse-case, level, nonclaim and audit data of
    Definitions~\ref{def:threshold-attachment} and~\ref{def:role-projection}.
\end{itemize}

It does not assert:
\begin{itemize}[topsep=2pt,itemsep=1pt]
  \item decomposition uniqueness or canonicity;
  \item six active projections, or six active nontrivial witnesses, in
    every description;
  \item scope-free residual exhaustion, or any claim about features
    outside the covered scope;
  \item that every finite audited feature of a description in \FATCD{}
    satisfies (R1)--(R5); Remark~\ref{rem:recognition-open} exhibits a
    feature with label (x) and a feature with label (xi);
  \item unrestricted exact-six, or coverage of all emergence
    phenomena;
  \item that broader domains require no further work.
\end{itemize}

\subsection{Global nonclaims}\label{subsec:global-nonclaims}

Each nonclaim below is preserved throughout the paper and is enforced
locally by the theorem-chain result that would otherwise overread it.

\begin{itemize}[topsep=2pt,itemsep=1pt]
  \item \emph{No full six-primitives algebra.} Scoped exact-six is
    proved over \FATCD{}; it is not claimed that \Pone{}--\Psix{}
    generate a global algebra of closure operations
    (Remark~\ref{rem:atlas-global-nonclaims},
    Remark~\ref{prop:overclaim-stress}).
  \item \emph{No global primitive-independence theorem.} Typed
    non-collapse (Proposition~\ref{prop:typed-non-collapse}) is the
    only claim made at that level; no global independence notion is defined
    or claimed.
  \item \emph{No empirical realization.} Empirical-bridge admissibility
    (Definition~\ref{prop:empirical-bridge}) admits
    threshold-dependent, level-indexed, instrument-visible substrate
    evidence; it does not assert that any \(D \in \FATCD{}\) is itself
    a substrate observation.
  \item \emph{No unsupported unique selected lift.} The displayed atlas lifts
    are selected and, except where uniqueness is proved and recorded (as for
    a descended map), non-unique; a local uniqueness assertion is admissible
    only with a recorded proof (Definition~\ref{prop:forgetting-lifts}).
  \item \emph{No unsupported lower-to-higher determinacy.} Behavioral or
    verification content is not claimed to determine structural-categorical
    content without a recorded proof, as for a descended map.
  \item \emph{No instrument-free formulation.} Admissible visibility claims carry
    an instrument-relative visibility record
    (Definition~\ref{prop:visibility}).
  \item \emph{No internal instrument-family completeness.} The
    accepted instrument family is not asserted to be sufficient for
    every exact-six judgment formulable in the same regime.
  \item \emph{No unrestricted all-emergence coverage.}
    Theorem~\ref{thm:scoped-exact-six} is scoped to \FATCD{}.
  \item \emph{No uniqueness of primitive presentations.} Distinct
    admissible presentations of the same role may coexist.
  \item \emph{Cross-level comparison requires translation.} Direct
    comparison of claims at different levels is inadmissible without a
    translation record (Definition~\ref{prop:level-profile}).
\end{itemize}

\subsection{The endogenous-instrument meta-limit}\label{subsec:meta-limit}

An adjacent meta-theoretic question, deferred to future work, is
whether the accepted internal instrument family can itself certify
both its soundness and its completeness for exact-six judgments
formulable inside the same regime. The notions of an internal argument and of
the soundness and completeness of an instrument family are not defined in this
paper; what follows is a heuristic target, not a mathematical statement. A
plausible sharp formulation is
that no internal argument should be expected to accomplish all three
of the following at once:
\begin{enumerate}[topsep=2pt,itemsep=1pt]
  \item establish the scoped exact-six claim;
  \item certify the soundness of the accepted instrument family;
  \item certify the completeness of that family for exact-six
    judgments expressible inside the same regime.
\end{enumerate}
Such self-certification limits are of a piece with the classical
G{\"o}del--L{\"o}b tradition on internal consistency and provability
statements~\citep{Godel1931,Lob1955}. This question is not part of
the proof of Theorem~\ref{thm:scoped-exact-six} and does not weaken
or retract the theorem's scoped conclusion. It is taken up as a
follow-on direction in Section~\ref{sec:future-work}.
